\documentclass[11pt]{article}
\usepackage[a4paper,margin=25mm]{geometry}
\usepackage{amsmath,amssymb,amsthm,mathtools,hyperref}
\newtheorem{theorem}{Theorem}\newtheorem{lemma}[theorem]{Lemma}
\newcommand{\dd}{\,d}\newcommand{\R}{\mathbb R}
\title{The original equilibrium density, its second moment, and a signed Gaussian return}
\author{Split-Zero research programme: independent derivation}
\date{20 September 2026}
\begin{document}\maketitle
This note proves the equilibrium and finite-constant part of the incoming
Gaussian continuation. It does not assume a spectral convergence theorem
for the arithmetic quotient. In particular, the positivity of the measure
below is proved by an explicit integral comparison; it is not inferred
from an unproved spectral limit. The potential, endpoint coordinates and
probability masses are those of EIQ1--33 in the programme source
\cite{EIQ}. Legendre's elliptic-integral conventions are those of Carlson
\cite{Carlson}; the original equation TeX was read for the definitions.
All identities and estimates used in this note are proved below.

\section{Endpoints, density and equilibrium}
For $0<s\le1$, set
\begin{equation}\tag{GM1}
 \begin{gathered}
 V_s(x)=\frac\pi s\sqrt x-\frac2s\log x\quad(x>0),\qquad
 \\
 K(\kappa)=\int_0^{\pi/2}\frac{\dd\theta}{\sqrt{1-\kappa^2\sin^2\theta}},
 \quad E(\kappa)=\int_0^{\pi/2}\sqrt{1-\kappa^2\sin^2\theta}\dd\theta.
 \end{gathered}
\end{equation}
The definitions of $K,E$ agree exactly with \cite[19.2.4, 19.2.5,
19.2.8(a,b)]{Carlson}. Put $r=\sqrt{1-\kappa^2}$. There is exactly one
$r\in(0,1)$ for which
\begin{equation}\tag{GM2}
 \frac{rK}{E}=\frac1{1+s},\qquad
 u=\frac2K,\quad v=\frac{2(1+s)}E,\quad r=\frac uv,
 \quad a=u^2,\quad b=v^2.
\end{equation}
Indeed, differentiation under the defining integrals and integration of
$d(\sin\theta\cos\theta/\sqrt{1-\kappa^2\sin^2\theta})$ give
$E_\kappa'=(E-K)/\kappa$ and
$K_\kappa'=E/[\kappa(1-\kappa^2)]-K/\kappa$.
They imply
\begin{equation}\tag{GM3}
 \frac{d}{dr}\frac{rK}E
 =\frac{(K-E)(E-r^2K)}{(1-r^2)E^2}>0.
\end{equation}
Both factors are positive by the defining integrals. Dominated convergence
gives $rK/E\to0$ at $r\downarrow0$ and $rK/E\to1$ at $r\uparrow1$:
for the first limit the integrand of $rK$ is bounded by one and tends
to zero except at the last endpoint, while $E\to1$.
This proves existence, uniqueness and continuous dependence.

With $R_w=\sqrt{(a+w)(b+w)}$ define
\begin{equation}\tag{GM4}
 \rho_s(x)=\frac{\sqrt{(b-x)(x-a)}}{4\pi s x}
       \int_0^\infty\frac{\sqrt w}{(x+w)R_w}\dd w\quad(a<x<b),
 \qquad \rho_s(x)=0\quad(x\notin(a,b)).
\end{equation}
Here is a derivation verifying its probability mass and the original
equilibrium problem. It also specifies the transform needed for its moment.
Write $\alpha=2/s$, $\beta=\pi/s$, $m=(a+b)/2$, $R_0=uv$ and
$R(z)=\sqrt{(z-a)(z-b)}$, with $R(z)/z\to1$ at infinity.
The substitution $x=a\cos^2\theta+b\sin^2\theta$ gives
\begin{equation}\tag{GM5}
 \int_a^b\frac{V_s'(x)}{\sqrt{(b-x)(x-a)}}\dd x=0,
 \qquad
 \int_a^b\frac{xV_s'(x)}{\sqrt{(b-x)(x-a)}}\dd x=2\pi.
\end{equation}
The required elementary integrals, with numerators $x^{-1/2}$,
$x^{-1}$ and $x^{1/2}$, are $2K/v$, $\pi/(uv)$ and $2vE$.
The substitution $w=xt^2$ gives
$\int_0^\infty[\sqrt w(x+w)]^{-1}\dd w=\pi/\sqrt x$.
Tonelli's theorem and (GM5) now imply
\begin{equation}\tag{GM6}
 \frac\alpha{R_0}=\frac\beta{2\pi}
       \int_0^\infty\frac{\dd w}{\sqrt wR_w},\qquad
 \int_0^\infty\frac{1-w/R_w}{\sqrt w}\dd w=2vE.
\end{equation}
Thus $h(x):=2\pi\rho_s(x)/\sqrt{(b-x)(x-a)}$ has the two forms
\begin{equation}\tag{GM7}
 h(x)=\frac\alpha{xR_0}-\frac\beta{2\pi}
       \int_0^\infty\frac{\dd w}{\sqrt w(x+w)R_w}
 =\frac\beta{2\pi x}\int_0^\infty\frac{\sqrt w}{(x+w)R_w}\dd w>0.
\end{equation}
The arcsine transform, proved by $x=m+(b-a)\cos\theta/2$ followed
by $t=\tan(\theta/2)$, is
$\pi^{-1}\int_a^b[(z-x)\sqrt{(b-x)(x-a)}]^{-1}\dd x=1/R(z)$.
Polynomial division and partial fractions consequently give
\begin{equation}\tag{GM8}
 \frac1\pi\int_a^b\frac{\sqrt{(b-x)(x-a)}}{x+w}\dd x=m+w-R_w,
 \quad
 \frac1\pi\int_a^b\frac{\sqrt{(b-x)(x-a)}}{(x+w)(z-x)}\dd x
 =1-\frac{R_w+R(z)}{z+w}=:H_w(z).
\end{equation}
Integrating (GM7) using (GM8), and then (GM6), proves
\begin{equation}\tag{GM9}
 \int_a^b\rho_s(x)\dd x
 =\frac\alpha{2R_0}(m-R_0)-\frac\beta{4\pi}
  \int_0^\infty\frac{(m+w)/R_w-1}{\sqrt w}\dd w
 =-\frac\alpha2+\frac{\beta vE}{2\pi}=1.
\end{equation}
All auxiliary integrals converge at zero and infinity. The Cauchy transform
has the exact expression
\begin{equation}\tag{GM10}
 G(z)=\int_a^b\frac{\rho_s(x)}{z-x}\dd x
 =\frac\alpha{2R_0}H_0(z)-\frac\beta{4\pi}
       \int_0^\infty\frac{H_w(z)}{\sqrt wR_w}\dd w
 =\frac{V_s'(z)-R(z)h(z)}2.
\end{equation}
The last equality follows by substituting (GM8); (GM6) cancels the
constant integral and the elementary Stieltjes integral supplies
$\beta/(4\sqrt z)$. It holds first outside both $[a,b]$ and the
negative real axis, and then throughout the domain of the Cauchy transform.

For completeness, this density is the unique equilibrium probability
for $\int V_s\dd\mu-\iint\log|x-y|\dd\mu\dd\mu$ on $(0,\infty)$.
The effective potential
$U(x)=V_s(x)-2\int\log|x-t|\rho_s(t)\dd t$ is continuous.
The boundary values of (GM10) give $U'=0$ on $(a,b)$ and
$U'=R(x)h(x)<0$ on $(0,a)$, $U'>0$ on $(b,\infty)$.
Hence $U=\ell$ on $[a,b]$ and $U>\ell$ elsewhere.
To justify the minimizing conclusion, let $\mu$ have finite energy,
$\nu=\mu-\rho_s\dd x$, and observe that the confining potential
controls $\sqrt x+|\log x|$. All cross logarithmic integrals are therefore
absolutely integrable; the negative logarithmic singularity in the cross
term is bounded after integration against the bounded compact density.
The scalar integral
$-\log d=\tfrac12\int_0^\infty(e^{-td^2}-e^{-t})\dd t/t$
and Gaussian Fourier integration give
\[
 -\iint\log|x-y|\dd\nu(x)\dd\nu(y)
 =\frac1{2\sqrt{4\pi}}\int_0^\infty t^{-3/2}
       \int_\R e^{-\xi^2/(4t)}|\widehat\nu(\xi)|^2\dd\xi\dd t\ge0.
\]
Truncation of the scalar integral is bounded by $|\log d|$, so dominated
convergence justifies this identity. Equality forces $\widehat\nu=0$ and
then $\nu=0$, by convolution with Gaussian approximate identities.
The energy difference is the sum of this nonnegative quantity and
$\int(U-\ell)\dd\mu$. It is strictly positive unless $\mu=\rho_s\dd x$.
Infinite energy cannot improve the finite value. This proves the assertion.

\section{Exact first moment in the squared coordinate}
An equivalent Cauchy-transform formula is
\begin{equation}\tag{GM11}
 G(z)=\frac{R(z)}{2\pi}
 \int_a^b\frac{V_s'(x)}{(z-x)\sqrt{(b-x)(x-a)}}\dd x.
\end{equation}
To verify it, subtract the right side from (GM10). Its boundary jumps
vanish: the right side has boundary sum $V_s'$ and the same jump as
the Cauchy transform, as follows either from (GM7) and the Stieltjes
integral or directly by the partial fractions in (GM8).
The difference extends across the endpoints and vanishes at infinity
by (GM5), so is zero. Equivalently, inserting
$V_s'=\beta/(2\sqrt x)-\alpha/x$ and the Stieltjes integral in (GM11)
reduces it term by term to (GM10).
Let $I_j=\int_a^b x^jV_s'(x)/\sqrt{(b-x)(x-a)}\dd x$.
Since $I_0=0$, $I_1=2\pi$ and $R(z)=z-m+O(z^{-1})$, expansion gives
\begin{equation}\tag{GM12}
 \int x\rho_s(x)\dd x=\frac{I_2}{2\pi}-m
 =\frac{\beta v^3}{2\pi}J_3-\left(\frac\alpha2+1\right)m,
 \quad J_3=\int_0^{\pi/2}(1-\kappa^2\sin^2\theta)^{3/2}\dd\theta.
\end{equation}
The exact reduction is
\begin{equation}\tag{GM13}
 3J_3=2(2-\kappa^2)E-(1-\kappa^2)K.
\end{equation}
Indeed the integrand of the left side minus the right side equals
$\kappa^2\,d(\sin\theta\cos\theta
\sqrt{1-\kappa^2\sin^2\theta})/d\theta$; its endpoint values vanish.
Substitute $\alpha=2/s$, $\beta=\pi/s$, $uK=2$ and $vE=2(1+s)$
in (GM12)--(GM13). This proves the requested original-coordinate value
\begin{equation}\tag{GM14}
 \boxed{\displaystyle\int x\rho_s(x)\dd x
 =\frac{(1+s)(u^2+v^2)-2uv}{6s}.}
\end{equation}

\section{A positive probability after the actual relation subtraction}
Define $h_0(y)=(2\pi)^{-1}\operatorname{arcosh}(2/|y|)$ for $0<|y|<2$,
and zero elsewhere. Its integrable singularity at zero is understood in
the measure sense. It is the mixture, for $0<t<1$, of the arcsine
probabilities on $[-2t,2t]$:
\begin{equation}\tag{GM15}
 h_0(y)=\int_{|y|/2}^1\frac{\dd t}{\pi\sqrt{4t^2-y^2}},\qquad
 \int y^{2j}h_0(y)\dd y=\frac{\binom{2j}j}{2j+1}.
\end{equation}
The second identity follows by $y=2t\cos\theta$ and integration over
$t\in[0,1]$, using
$\pi^{-1}\int_0^\pi\cos^{2j}\theta\dd\theta=\binom{2j}j/4^j$.
This last formula follows successively by integration by parts, starting
from one at $j=0$.

For $0<s\le1$ define, in the original $y$ coordinate,
\begin{equation}\tag{GM16}
 \dd\nu_s(y)=\left[h_0\left(\frac y{1+s}\right)
                    -s|y|\rho_s(y^2)\right]\dd y.
\end{equation}
This is a positive probability supported in $[-2(1+s),2(1+s)]$.
Here is a direct proof, independent of any arithmetic spectral limit.
For $a<x<b$, (GM4) gives
\begin{align}\tag{GM17}
 s\sqrt x\rho_s(x)
 &\le\frac{\sqrt{b-x}}{4\pi}
          \int_0^\infty\frac{\dd w}{(x+w)\sqrt{b+w}}\notag\\
 &=\frac1{2\pi}\operatorname{arcosh}\frac v{\sqrt x}
 \le\frac1{2\pi}\operatorname{arcosh}\frac{2(1+s)}{\sqrt x}.\notag
\end{align}
The first inequality uses separately $\sqrt{x-a}/\sqrt x\le1$ and
$\sqrt w/\sqrt{a+w}\le1$. The substitution $t=\sqrt{b+w}$ evaluates
the integral as
$2\operatorname{arcosh}(v/\sqrt x)/\sqrt{b-x}$.
The last inequality uses $E\ge1$, hence $v\le2(1+s)$.
Outside this interval the subtracted density vanishes, so (GM17)
proves positivity everywhere. The pushforward of
$|y|\rho_s(y^2)\dd y$ under $y\mapsto y^2$ is exactly
$\rho_s(x)\dd x$, with both signs of $y$ included. Its mass is one.
The first term of (GM16) has mass $1+s$, so $\nu_s$ has mass one.
At $s=0$ put $\dd\nu_0=h_0(y)\dd y$.

Its second moment is consequently
\begin{equation}\tag{GM18}
 \boxed{\mathfrak m(s)=\int y^2\dd\nu_s(y)
 =\frac23(1+s)^3-\frac{(1+s)(u^2+v^2)-2uv}{6}.}
\end{equation}
In particular $\mathfrak m(0)=2/3$ and
$\mathfrak m(1)=(16-u^2-v^2+uv)/3$.
For the endpoint limit $s\downarrow0$, (GM2) gives
$r\uparrow1$ and $u,v\to4/\pi$; the formula agrees continuously
with this definition at zero.

\section{Finite lower bounds, with no fitted elliptic constants}
First, $r(1)>1/8$. The elementary tangent substitution gives
\[
 K=\int_0^\infty\frac{\dd t}{\sqrt{(1+t^2)(1+r^2t^2)}}
 \le2\operatorname{arsinh}(1)+\log(1/r).
\]
For this bound split the integral at $1$ and $1/r$, use the change
$t=1/(rw)$ in the final part, and bound the middle integrand by $1/t$.
At $r=1/8$ this is less than $39/10<4$: the elementary inequalities
$\operatorname{arsinh}(1)<9/10$ and $\log2<7/10$ follow from
$\sqrt2<99/70$ and the positive Taylor sums through degree five for
$e^{9/10}$ and $e^{7/10}$ respectively.
Thus $rK/E<1/2$ at $r=1/8$ and (GM3) proves the assertion.
Also $K>E$ gives $r(1)<1/2$.

If $r=r(1)\ge3/20$, the decreasing function $1-r+r^2$ on $(0,1/2)$
and $v<4$ give $v^2(1-r+r^2)<349/25$.
If $1/8<r\le3/20$, differentiation under the integral in (GM1) gives
\begin{equation}\tag{GM19}
 E(r)-1=\int_0^r t\int_0^{\pi/2}
 \frac{\sin^2\theta}{\sqrt{\cos^2\theta+t^2\sin^2\theta}}
 \dd\theta\dd t
 \ge\frac{r^2(K-E)}{2(1-r^2)}.
\end{equation}
Use $K=E/(2r)$ and
$r(1-2r)/[4(1-r^2)]\ge7/320$ on this interval.
It follows that $E\ge320/313$ and
\[
 v^2(1-r+r^2)\le\frac{57}{4}\left(\frac{313}{320}\right)^2
 <\frac{349}{25}.
\]
Both cases prove $\mathfrak m(1)>17/25$ with exact rational constants.

The same argument proves the claimed uniform bound. For $0<s\le1$,
(GM3) gives $r(s)\ge r(1)>1/8$; $K$ decreases with $r$, so $K<4$
and $u>1/2$. Furthermore $v\ge u$, $v\le2(1+s)$, and
$(1+s)u/v=E/K<1$. With $T=1+s$, the exact identity
\begin{equation}\tag{GM20}
 6\mathfrak m(s)-uv=T(4T^2-v^2)+u(v-Tu)\ge0
\end{equation}
therefore gives $\mathfrak m(s)\ge uv/6>1/24$. The assertion is also
true at $s=0$, where its value is $2/3$.

\section{An explicit finite series certificate strengthening the endpoint}
The accompanying exact/radius computation strengthens the finite endpoint
bound to
\begin{equation}\tag{GM21}
 \frac{16010167095}{10^{11}}<r(1)<\frac{16010167096}{10^{11}},
 \qquad \boxed{\mathfrak m(1)>\frac{10249}{10000}.}
\end{equation}
Here are the full mathematical error bounds used by the certificate.
For $z=1-r^2\in(0,1)$ and
$a_n=(\binom{2n}n/4^n)^2$, expansion of the defining integrands yields
\begin{align}\tag{GM22}
 \frac{2K}{\pi}&=\sum_{n=0}^{J}a_nz^n+R_K,
 &0<R_K\le\frac{z^{J+1}}{1-z},\notag\\
 \frac{2E}{\pi}&=1-\sum_{n=1}^{J}\frac{a_nz^n}{2n-1}-R_E,
 &0<R_E\le\frac{z^{J+1}}{(2J+1)(1-z)}.\notag
\end{align}
For $K$ use the binomial series and the cosine moment in (GM15).
For $E$ the binomial coefficient is the negative of the corresponding
coefficient for $K$ divided by $2n-1$. Termwise integration is justified
by uniform convergence because $z<1$. The tail bounds use $0<a_n\le1$.
The recurrence $a_0=1$, $a_{n+1}=a_n((2n+1)/(2n+2))^2$ retains every
coefficient. The certificate uses $J=2048$ and outward-rounded real
interval arithmetic with 240-bit midpoints. It proves that $rK/E-1/2$
is negative at the lower rational in (GM21), positive at the upper,
and then evaluates (GM18) on the entire resulting interval.
The ratio test is independent of $\pi$. The endpoints $u=2/K$,
$v=4/E$ use a rigorous $\pi$ interval obtained from
\[
 \pi=16\arctan(1/5)-4\arctan(1/239).
\]
For completeness, the tangent addition formula gives
$\tan(4\arctan(1/5)-\arctan(1/239))=1$ exactly;
the angle is in $(0,\pi/2)$, proving the identity.
Alternating arctangent sums with their first omitted terms bound both
arctangents (respectively 24 and 8 terms). Thus no floating point
approximation to an unbounded remainder is used. All rational endpoints,
interval enclosures and pass/fail comparisons are saved in the check
receipt. The general bound (GM20) remains an analytic proof over the
entire continuum; sampled parameter checks are not substituted for it.

\section{The complete signed Gaussian expression}
Let
\begin{equation}\tag{GM23}
 H_s(t)=\int e^{-ty^2}\dd\nu_s(y),\qquad
 \mathfrak A(t)=2[-H_0'(t)+H_1'(t)]\quad(t\ge0).
\end{equation}
All these derivatives exist by compact support. If
$I_0(z)=\pi^{-1}\int_0^\pi e^{z\cos\theta}\dd\theta$, the arcsine
mixture (GM15) gives exactly
\begin{equation}\tag{GM24}
 H_0(t)=\int_0^1e^{-2tw^2}I_0(2tw^2)\dd w,\qquad
 H_s(t)=(1+s)H_0((1+s)^2t)-s\int e^{-tx}\rho_s(x)\dd x.
\end{equation}
The second identity is the two-to-one pushforward $x=y^2$ of (GM16),
including its full mass. Since $\nu_0,\nu_1$ are positive probabilities
with squared supports in $[0,4]$ and $[0,16]$, respectively,
\[
 -H_0'(1/1024)\le2/3,\qquad
 -H_1'(1/1024)\ge\frac{63}{64}\mathfrak m(1).
\]
Here $e^{-x/1024}\ge1-x/1024\ge63/64$ on the second support.
The completely analytic bound $\mathfrak m(1)>17/25$ proves the incoming
$\mathfrak A(1/1024)<-13/2400$. The finite-series certificate strengthens
this to
\begin{equation}\tag{GM25}
 \boxed{\mathfrak A(1/1024)
 <\frac43-\frac{63}{32}\frac{10249}{10000}
 =-\frac{657061}{960000}.}
\end{equation}
This statement is about the explicitly constructed equilibrium receiver.
Its transfer to the original finite arithmetic observation matrices
requires the actual spectral convergence argument. This note proves
neither that convergence nor the individual projected phase signs.

\begin{thebibliography}{9}
\bibitem{Carlson} B. C. Carlson, \emph{Elliptic Integrals}, Chapter 19 of
the NIST Digital Library of Mathematical Functions, version 1.2.8,
15 September 2026, equations 19.2.4, 19.2.5 and 19.2.8(a,b).
\url{https://dlmf.nist.gov/19.2}. Carlson credits L. M. Milne-Thomson's
Chapter 17 in Abramowitz--Stegun for part of the Legendre-integral material.
The four original equation TeX files are retained in \texttt{author\_sources}.
\bibitem{EIQ} Split-Zero research programme, \emph{Equilibrium measure for
the original Gamma relation potential}, EIQ1--33, cumulative source
\texttt{09\_UPDATED\_JOINT\_NOTE.tex}, original potential, elliptic endpoints,
density, Cauchy transform and energy proof. The private reading ledger
records the exact source hash and coverage; this note rederives every
equilibrium identity needed here.
\end{thebibliography}
\end{document}
